\Article{35}
The following problem deserves a separate treatment, both because of its practical usefulness and because of the elegance of its solution.

\textit{To find the changes in the most plausible values of the unknowns produced by the addition of a new equation, and the weights of the new determinations.}

We shall retain the notation used above, so that the original equations, reduced to weight \(=1\), are these: \(v=0,\allowbreak v'=0,\allowbreak v''=0\) etc.; for the indefinite su\SourceLeaf{048}%
m we write \(vv+v'v'+v''v''\quad\text{etc.}=\Omega\); and \(\xi,\eta,\zeta\) etc. are the partial differential quotients
\[
\frac{d\Omega}{2dx},\quad\frac{d\Omega}{2dy},\quad\frac{d\Omega}{2dz}\quad\text{etc.}
\]
and finally indefinite elimination gives
\[
\left.\begin{aligned}
x&=A+[\alpha\alpha]\xi+[\alpha\beta]\eta+[\alpha\gamma]\zeta+\text{etc.}\\
y&=B+[\alpha\beta]\xi+[\beta\beta]\eta+[\beta\gamma]\zeta+\text{etc.}\\
z&=C+[\alpha\gamma]\xi+[\beta\gamma]\eta+[\gamma\gamma]\zeta+\text{etc.}
\end{aligned}\right\}\ (\mathrm{I})
\]
Now let us suppose that a new equation \(v^*=0\) is added (approximately true, and with weight \(=1\)), and let us investigate how much this will change both the most plausible values of the unknowns \(A,B,C\) etc. and the coefficients \([\alpha\alpha],[\alpha\beta]\) etc.

Let us put \(\Omega+v^*v^*=\Omega^*,\allowbreak\frac{d\Omega^*}{2dx}=\xi^*,\allowbreak\frac{d\Omega^*}{2dy}=\eta^*,\allowbreak\frac{d\Omega^*}{2dz}=\zeta^*\) etc., and suppose that elimination then gives
\[
x=A^*+[\alpha\alpha^*]\xi^*+[\alpha\beta^*]\eta^*+[\alpha\gamma^*]\zeta^*\quad\text{etc.}
\]
Finally, let
\[
v^*=fx+gy+hz+\text{etc.}+k
\]
so that, on substituting for \(x,y,z\) etc. their values from (I), we obtain
\[
v^*=F\xi+G\eta+H\zeta+\text{etc.}+K
\]
and let us put
\[
Ff+Gg+Hh+\text{etc.}=\omega
\]
Clearly \(K\) will be the most plausible value of the function \(v^*\) as far as it follows from the original equations, without regard to the value \(0\) supplied by the additional observation, and \(\frac1\omega\) will be the weight of this determination.

We now have
\[
\xi^*=\xi+fv^*,\quad\eta^*=\eta+gv^*,\quad\zeta^*=\zeta+hv^*\quad\text{etc.}
\]
and therefore
\[
\begin{aligned}
&F\xi^*+G\eta^*+H\zeta^*+\text{etc.}+K\\
&\qquad=v^*(1+Ff+Gg+Hh+\text{etc.})
\end{aligned}
\]
or
\[
v^*=\frac{F\xi^*+G\eta^*+H\zeta^*+\text{etc.}+K}{1+\omega}
\]
\SourceLeaf{049}
\noindent Similarly we obtain
\[
\begin{aligned}
x&=A+[\alpha\alpha]\xi^*+[\alpha\beta]\eta^*+[\alpha\gamma]\zeta^*+\text{etc.}\\
&\qquad-v^*(f[\alpha\alpha]+g[\alpha\beta]+h[\alpha\gamma]+\text{etc.})\\
&=A+[\alpha\alpha]\xi^*+[\alpha\beta]\eta^*+[\alpha\gamma]\zeta^*+\text{etc.}-Fv^*\\
&=A+[\alpha\alpha]\xi^*+[\alpha\beta]\eta^*+[\alpha\gamma]\zeta^*+\text{etc.}\\
&\qquad-\frac{F}{1+\omega}(F\xi^*+G\eta^*+H\zeta^*+\text{etc.}+K)
\end{aligned}
\]
From this we therefore conclude that
\[
A^*=A-\frac{FK}{1+\omega}
\]
which will be the most plausible value of \(x\) from \textit{all} the observations;
\[
[\alpha\alpha^*]=[\alpha\alpha]-\frac{FF}{1+\omega}
\]
and the weight of this determination is therefore
\[
=\frac1{[\alpha\alpha]-\dfrac{FF}{1+\omega}}
\]
In exactly the same way we find the most plausible value of \(y\), based on \textit{all} the observations,
\[
B^*=B-\frac{GK}{1+\omega}
\]
and the weight of this determination is
\[
=\frac1{[\beta\beta]-\dfrac{GG}{1+\omega}}
\]
and so on. Q.\ E.\ I.

Let us add a few remarks to this solution.

I. On substituting these new values \(A^*,B^*,C^*\) etc., the function \(v^*\) acquires the most plausible value
\[
K-\frac{K}{1+\omega}(Ff+Gg+Hh+\text{etc.})=\frac{K}{1+\omega}
\]
And since we have indefinitely
\[
v^*=\frac{F}{1+\omega}.\xi^*+\frac{G}{1+\omega}.\eta^*+\frac{H}{1+\omega}.\zeta^*+\text{etc.}+\frac{K}{1+\omega}
\]
\SourceLeaf{050}%
the weight of this determination is found by the principles of art.\ 29 to be
\[
=\frac{1+\omega}{Ff+Gg+Hh+\text{etc.}}=\frac1\omega+1
\]
The same results follow directly by applying the rule given at the end of art.\ 21; namely, the system of original equations had supplied the determination \(v^*=K\) with weight \(=\frac1\omega\), and the new observation then gave another determination, independent of the former, \(v^*=0\), with weight \(=1\); combining them gives the determination \(v^*=\frac{K}{1+\omega}\) with weight \(=\frac1\omega+1\).

II. It further follows that, since for \(x=A^*,\allowbreak y=B^*,\allowbreak z=C^*\) etc. we must have \(\xi^*=0,\allowbreak\eta^*=0,\allowbreak\zeta^*=0\) etc., the same values give
\[
\xi=-\frac{fK}{1+\omega},\quad\eta=-\frac{gK}{1+\omega},\quad\zeta=-\frac{hK}{1+\omega}\quad\text{etc.}
\]
and also, since indefinitely
\(\Omega=\xi(x-A)+\eta(y-B)+\zeta(z-C)+\text{etc.}+M\),
\[
\begin{aligned}
\Omega&=\frac{KK}{(1+\omega)^2}(Ff+Gg+Hh+\text{etc.})+M\\
&=M+\frac{\omega KK}{(1+\omega)^2}
\end{aligned}
\]
and finally, since indefinitely
\(\Omega^*=\Omega+v^*v^*\),
\[
\Omega^*=M+\frac{\omega KK}{(1+\omega)^2}+\frac{KK}{(1+\omega)^2}=M+\frac{KK}{1+\omega}
\]
III. Comparing this with what we taught in art.\ 30, we observe that the function \(\Omega\) here attains the smallest value it can take for the specified value of the function \(v^*=\frac{K}{1+\omega}\).

\Article{36}
For another problem, related to the preceding one, namely

\textit{To investigate the changes in the most plausible values of the unknowns arising from a change in the weight of one of the original observations, and the weights of the new determinations}

\noindent we shall give only the solution here, omitting for brevity the proof, which is easily completed on the model of the preceding article.

Let us suppose that only after completing the calculation we notice that one of the observations has been assigned a weight either too small or too large; e.\ g. that the first observation, which gave \(V=L\), should be assigned, instead of the weight \(p\) used in the calculation, more correctly a weight \(=p^*\). There will then be no need to repeat the entire calculation; instead, the corrections can be computed more conveniently by the following formulas.

\SourceLeaf{051}%
The corrected most plausible values of the unknowns will be these:
\[
\begin{aligned}
x&=A-\frac{(p^*-p)\alpha\lambda}{p+(p^*-p)(a\alpha+b\beta+c\gamma+\text{etc.})}\\[7pt]
y&=B-\frac{(p^*-p)\beta\lambda}{p+(p^*-p)(a\alpha+b\beta+c\gamma+\text{etc.})}\\[7pt]
z&=C-\frac{(p^*-p)\gamma\lambda}{p+(p^*-p)(a\alpha+b\beta+c\gamma+\text{etc.})}
\end{aligned}
\]
etc., and the weights of these determinations will be found by dividing unity respectively by
\[
\begin{aligned}
&[\alpha\alpha]-\frac{(p^*-p)\alpha\alpha}{p+(p^*-p)(a\alpha+b\beta+c\gamma+\text{etc.})}\\[7pt]
&[\beta\beta]-\frac{(p^*-p)\beta\beta}{p+(p^*-p)(a\alpha+b\beta+c\gamma+\text{etc.})}\\[7pt]
&[\gamma\gamma]-\frac{(p^*-p)\gamma\gamma}{p+(p^*-p)(a\alpha+b\beta+c\gamma+\text{etc.})}\quad\text{etc.}
\end{aligned}
\]
This solution also includes the case where, after completing the calculation, we realize that one of the observations should have been rejected altogether, since this is the same as putting \(p^*=0\); and similarly the value \(p^*=\infty\) refers to the case where the equation \(V=L\), which was treated as approximate in the calculation, actually has absolute precision.

Moreover, whenever either \textit{several} new equations are added to those on which the calculation was based, or we realize that \textit{several} of them have been assigned incorrect weights, the calculation of the corrections would become too complicated; in such a case it will therefore be better to redo the calculation from the beginning.

\clearpage
